\documentclass{article}
\usepackage[left=2cm, right=2cm, top=2cm, bottom=2cm]{geometry}
\usepackage{graphicx}
\usepackage{amsmath}
\usepackage{amssymb}
\usepackage{amsthm}
\usepackage{fancyhdr}
\usepackage{verbatim}
\usepackage{listings}
\usepackage{xcolor}
\usepackage{pdfpages}
\usepackage{cancel}
\usepackage{physics}

\lstset{
    language=C++,
    basicstyle=\ttfamily\footnotesize,
    keywordstyle=\color{blue},
    commentstyle=\color{green},
    stringstyle=\color{red},
    numbers=left,
    numberstyle=\tiny\color{gray},
    frame=single,
    breaklines=true,
    captionpos=b,
}


\title{MTH 553: Lecture \# 17}
\author{Cliff Sun}

\newtheorem{theorem}{Theorem}[section]
\newtheorem{lemma}[theorem]{Lemma}
\newtheorem{definition}[theorem]{Definition}
\newtheorem{conjecture}[theorem]{Conjecture}
\newtheorem{proposition}[theorem]{Proposition}
\newtheorem{corollary}[theorem]{Corollary}
\newtheorem{one minute paper}[theorem]{One Minute Paper}

\pagestyle{fancy}
\lhead{\textbf{\thepage}\ \ \nouppercase{\rightmark}}
\chead{MTH 553: Lecture \# 17}
\rhead{Cliff Sun}

\begin{document}

\maketitle

\section*{Lecture Span}
\begin{enumerate}
    \item Heat equation
\end{enumerate}

\section*{Eigenvalues of Laplacian}

Consider some bounded domain $\Omega$ in $\mathbb{R}^n$. Then 

\begin{align*}
    -\triangle u &= \lambda u \text{ in $\Omega$} \\
    u &= 0 \text{ on $\partial\Omega$}
\end{align*}

Here, $\lambda$ is the eigenvalue of the Laplacian for Dirichlet BC. And $u$ is the eigenfunction with $u \neq 0$. 

Goal: generalize fourier series to higher dimension on non-rectangular domains. 

\subsection*{Eigenfacts}

\begin{enumerate}
    \item $0 < \lambda_1 < \lambda_2 \leq \lambda_3 \leq \lambda_4 \leq \dots \rightarrow \infty$ with corresponding eigenfunctions $\phi_i$ such that $\{\phi_i\}$ are an orthonormal basis of $L^2(\Omega)$. 
    \begin{equation}
        \langle \phi_i, \phi_j \rangle = 0 \iff i \neq j
    \end{equation}
    and for all $f \in L^2(\Omega)$, we have that $$f = \sum_{j=1}^{\infty}a_j \phi_j$$ converging in $L^2$ with $a_j = \langle f, \phi_j \rangle_{L^2}$
    \item Positivity of fundamental mode (first eigenfunction) $\phi_1(x) > 0$ for all $\forall x \in \Omega$. (think of just one mode (parabola)) This means that $\phi_j$ must change sign for $j > 1$. 
    \item Pointwise convergence if $\partial \Omega$ is piecewise smooth, $f \in C^2\left( \overline{\Omega} \right)$ and if $f = 0$ on $\partial \Omega$, then $f = \sum_j a_j \phi_j$ converges pointwise. 
    In fact, uniformly and absolutely.
\end{enumerate}

\subsection*{Example $\Omega = (0, \pi) \subset \mathbb{R}^1$}

Here, 

\begin{equation}
    \phi_j(x) = \sqrt{\frac{2}{\pi}}\sin(jx)
\end{equation}

\subsection*{Example, rectangle in $R^2$ with lengths $a$ and $b$}

Eigenfunctions, $$\phi_{mn}(x,y) = a_{mn}\sin\left( \frac{m\pi x}{a} \right)\sin\left( \frac{n\pi x}{b} \right)$$

For $m,n \geq 1$. Indeed, $\phi_{mn}(x,y)$ is a eigenfunction of $\triangle$ with eigenvalues of 

\begin{equation}
    \lambda_{mn} = \left( \frac{m\pi}{a} \right)^2 + \left( \frac{n\pi}{a} \right)^2
\end{equation}

Take a square for example, then consider $a=b$. Therefore, the ordering of smallest to largest eigenvalue is 

\begin{enumerate}
    \item $\lambda_1 = \lambda_{11}$
    \item $\lambda_2 = \lambda_{21}$
    \item $\lambda_3 = \lambda_{12} = \lambda_{21}$
\end{enumerate}

There is a degeneracy. Other examples include balls, equilateral triangles, etc.

\subsection*{Applications}

Solving Poisson equation by eigenfunctions. Consider 

\begin{align*}
    \triangle u &= f \\
    u&= 0 \text{ on $\partial \Omega$}
\end{align*}

Now, write 

\begin{equation}
    f = \sum_k a_k \phi_k
\end{equation}

Then 

\begin{equation}
    u = \sum_k \left( -\frac{a_k}{\lambda_k} \right)\phi_k
\end{equation}

For the proof, take $\triangle$ on both sides and note that $-\triangle \phi_k = \lambda_k \phi_k$. Now, we can solve the wave equation using eigenfunctions. 

\subsection*{Wave equation with eigenfunctions}

\begin{align*}
    u_{tt} &= \triangle u \\
    u &= 0 
\end{align*}

Here, we seek a separated solution. Let 

\begin{equation}
    u(x,t) = \sum_k u_k(t)\phi_k(x)
\end{equation}

Then plug it into the wave equation:

\begin{equation}
    \sum_k u_k''(t)\phi_k(x) = \sum_k u_k(t)(-\lambda_k \phi_k(x))
\end{equation}

Then. we require that 

\begin{equation}
    u_k''(t) = \lambda_k u_k(t)
\end{equation}

Therefore, 

\begin{equation}
    u_k(t) = A_k\cos(\sqrt{\lambda_k}t) + B_k \sin(\sqrt{\lambda_k}t)
\end{equation}

If $u(x,0) = g(x)$, then 

\begin{equation}
    \sum_k A_k \phi_k(x) = g(x)
\end{equation}

Therefore, 

\begin{equation}
    A_k = \langle g, \phi_k \rangle_{L^2}
\end{equation}

Here, consider the term $\cos(\sqrt{\lambda_k} t)\phi_k(x)$. Then this is a standing wave with a mode shape $\phi_k(x)$ with a frequency of $$\frac{\sqrt{\lambda_k}}{2\pi}$$

\end{document}